\documentclass[11pt,a4paper]{article}
\usepackage[utf8]{inputenc}
\usepackage[T1]{fontenc}
\usepackage{geometry}
\usepackage{booktabs}
\usepackage{tabularx}
\usepackage{enumitem}
\usepackage{amsmath}
\usepackage{hyperref}

\geometry{margin=1in}

\title{\textbf{General Shape of NDVI Over Time Across Landscapes}}
\author{}
\date{}

\begin{document}

\maketitle

Over a full year, the general shape of a Normalized Difference Vegetation Index (NDVI) time-series curve for vegetated landscapes is a \textbf{bell-shaped (sinusoidal) seasonal wave}---starting low during dormancy, rising steeply during spring leaf-out (``green-up''), plateauing at peak summer growth, and declining during autumn senescence.

However, the exact curve shape changes distinctly depending on the \textbf{ecosystem and land management}.

\section{Common Landscape Curve Shapes}

\begin{table}[h!]
\centering
\renewcommand{\arraystretch}{1.35}
\begin{tabularx}{\textwidth}{>{\bfseries}p{3.8cm} p{3.8cm} X}
\toprule
Landscape Type & General Curve Shape & Key Characteristics \\
\midrule
Deciduous Forest \& Temperate Grassland & Smooth, single bell curve & Low baseline in winter ($0.1\text{--}0.2$), rapid spring green-up, high summer plateau ($0.7\text{--}0.85$), gradual autumn drop. \\
Single-Crop Agriculture & Sharp, narrow ``hat'' or trapezoid & Flat bare-soil baseline ($\sim 0.15$), very steep rise after planting, flat peak canopy, abrupt drop at harvest. \\
Double-Crop / Irrigated Agriculture & Bimodal (two distinct peaks) & Two separate bell curves within one calendar year (e.g., a winter wheat peak and a summer corn/soybean peak). \\
Evergreen / Tropical Forest & Flat high plateau & Consistently high ($0.75\text{--}0.90$) year-round with minimal seasonal dip, though tropical dry forests show a gentle wet/dry-season wave. \\
Semi-Arid / Mediterranean Shrubland & Asymmetric or pulse-driven wave & Tracks rainfall rather than temperature; spikes quickly after wet-season rains (e.g., winter/spring or summer monsoons) and tapers slowly during dry months. \\
Urban / Barren / Water & Flat low line & Remains low and relatively constant year-round ($< 0.0$ for water, $0.05\text{--}0.20$ for concrete/bare rock). \\
\bottomrule
\end{tabularx}
\caption{Comparison of annual NDVI time-series curve shapes across major land-cover types.}
\end{table}

\section{Key Phenological Stages on the Curve}

In remote sensing (land-surface phenology), analysts extract four key inflection points and metrics from this seasonal wave:

\begin{enumerate}[leftmargin=*]
    \item \textbf{Start of Season (SOS):} The inflection point where NDVI begins rising rapidly above the winter or dry-season baseline.
    \item \textbf{Peak of Season (POS):} The maximum NDVI value reached during peak biomass and canopy closure.
    \item \textbf{End of Season (EOS):} The point where NDVI drops back toward the dormant baseline.
    \item \textbf{Integrated NDVI (Area Under the Curve):} The total area between the seasonal curve and the baseline---directly analogous to summing $(\text{count} \times \text{bin value})$ over time---which serves as a proxy for \textbf{total annual primary productivity (biomass)}.
\end{enumerate}

\section{Real-World Noise in Raw Curves}

Raw satellite NDVI time series rarely look like a smooth mathematical curve on the first plot due to atmospheric and sensor limitations:

\begin{itemize}[leftmargin=*]
    \item \textbf{Cloud \& Snow Drops:} Clouds, haze, and snow artificially depress near-infrared (NIR) reflectance, causing sharp, sudden \textbf{downward spikes} that must be smoothed (e.g., using Savitzky--Golay filters or maximum-value composites).
    \item \textbf{Saturation Effect:} In very dense canopies (such as tropical forests or mature corn), NDVI flattens out near $0.85\text{--}0.90$ even as biomass continues to increase.
\end{itemize}

\end{document}
